% Cuerpo científico recuperado y distribuido en su residencia terminal.
\section{Clausura cofinal del criterio de Riemann--Weil}

La realización genealógica produce una familia cofinal y natural de formas
positivas cuya publicación coincide con el funcional completo de
Guinand--Weil. La extinción del terminal y la compatibilidad con las
inclusiones permiten pasar al límite sin introducir una lista de ceros ni el
signo buscado. Por el criterio reversible de Weil, todo cero no trivial
\(\rho\) de la función zeta de Riemann satisface
\[
  \operatorname{Re}\rho=\frac12.
\]
Los cálculos regulados de ventana finita son controles posteriores; no forman
parte de las premisas del cierre cofinal.

\paragraph{Factorización cofinal de la conclusión espectral.}
La cadena probatoria reúne, en el orden causal ya establecido, la forma
genealógica finita, su complemento positivo, la naturalidad bajo inclusiones y
la extinción compatible del término terminal. Cada ventana preserva los relojes
primos, el canal arquimediano, la orientación de las dos hojas y el registro de
frontera. La cofinalidad transporta conjuntamente esos datos al límite y la
identidad de Guinand--Weil identifica la forma resultante con su publicación
espectral clásica.

En consecuencia, la positividad no procede de una selección posterior de
ceros: desciende del operador positivo construido antes de la evaluación. La
involución funcional conserva la hoja espejo y fija la recta
\(\operatorname{Re}\rho=1/2\); la compatibilidad cofinal conserva esa fijación
en el paso desde las ventanas reguladas hasta la forma completa. Los cálculos
finitos cumplen así su función propia de verificación posterior, mientras que
la clausura matemática reside en la composición positiva y natural anterior.

